\chapter{The Integration Mode}\index{modes!Integration Mode|textbf}
\label{ch:integration}

\begin{versebox}
Those teachings that share a single root,\\
those proclaimed to share a single aim ---\\
integrate them all together:\\
everything converges to one flame.
\end{versebox}

\noindent We have now walked through fifteen modes of reading. Each one gave you a different lens: the Teaching Mode showed you the gratification-danger-escape pattern; the Investigation Mode divided the craving-side from the view-side; the Fitness Mode tested whether interpretations hold up; the Entry Mode showed you five doorways into every passage; the Requisites Mode mapped the chain of causes.

Now the Integration Mode says: put them all together\index{integration|textbf}. At a single basis-point, however many of these analyses converge, they should all be integrated --- woven into one understanding.\footnote{Pali: \textit{``Tattha katamo samāropano hāro? `Ye dhammā yaṃmūlā, ye cekatthā pakāsitā muninā\,''ti.''} The mnemonic fragment is ``those teachings that share a root, and those proclaimed by the Sage to have a single purpose'' (\textit{ye dhammā yaṃmūlā, ye cekatthā pakāsitā muninā}). The word \textit{samāropana} means ``mounting, placing together, integration'' --- the mode that takes the results of all the other modes and weaves them into a unified reading. \textit{``Ekasmiṃ padaṭṭhāne yattakāni padaṭṭhānāni otaranti, sabbāni tāni samāropayitabbāni.''} ``At a single basis-point, however many bases enter, all should be integrated.''}

The Netti identifies four types of integration: by basis, by synonym, by development, and by abandonment.

\section*{Integration by Basis}

Take the verse of all the Buddhas, which we have already met:

\begin{versebox}
Not to do any evil,\\
to cultivate the good,\\
to purify one's own mind ---\\
this is the teaching of all the Buddhas.
\end{versebox}

What is its basis? The three types of good conduct --- bodily, verbal, and mental. But now the Integration Mode traces the chain of bases further. Bodily and verbal good conduct are the aggregate of virtue\index{training aggregates (three)}. The mental good conduct of non-covetousness and non-ill-will is the aggregate of concentration. Right view is the aggregate of wisdom. Three aggregates.\footnote{Pali: \textit{``Tassa kiṃ padaṭṭhānaṃ? Tīṇi sucaritāni --- kāyasucaritaṃ vacīsucaritaṃ manosucaritaṃ --- idaṃ padaṭṭhānaṃ; tattha yaṃ kāyikañca vācasikañca sucaritaṃ, ayaṃ sīlakkhandho. Manosucarite yā anabhijjhā abyāpādo ca, ayaṃ samādhikkhandho. Yā sammādiṭṭhi, ayaṃ paññākkhandho.''} The Netti traces the basis-chain from the verse: good conduct → three aggregates (virtue, concentration, wisdom). The Integration Mode doesn't stop at the first basis but keeps going.}

Then: virtue and concentration together are calm. Wisdom is insight. Two practices.\footnote{Pali: \textit{``Sīlakkhandho ca samādhikkhandho ca samatho, paññākkhandho vipassanā.''}}

Then: the fruit of calm is liberation of the heart through the fading of lust. The fruit of insight is liberation through wisdom by the fading of ignorance. Two liberations.\footnote{Pali: \textit{``Samathassa phalaṃ rāgavirāgā cetovimutti, vipassanāya phalaṃ avijjāvirāgā paññāvimutti. Idaṃ padaṭṭhānaṃ.''} The chain of integration: verse → good conduct → three training-aggregates → calm and insight → two liberations. Five levels, each the basis of the next, all integrated into a single reading of a single verse.}

Five levels of basis --- good conduct, the three aggregates, calm and insight, and the two liberations --- all drawn from a single four-line verse. That is integration by basis.

The Netti adds another worked example: the Buddha's injunction to ``cut down both the forest and the underbrush.'' What is the forest? What is the underbrush? The Integration Mode traces four different basis-chains through this single metaphor. The forest is the five strands of sense-pleasure; the underbrush is craving. Or: the forest is grasping at the general sign (``a woman,'' ``a man''); the underbrush is grasping at specific features (``such eyes! such a voice!''). Or: the forest is the six internal and external sense-bases, uncomprehended; the underbrush is the fetters that arise from their meeting. Or: the forest is the underlying tendency; the underbrush is the active obsession.\footnote{Pali: \textit{``Vanaṃ vanathassa padaṭṭhānaṃ \ldots Vanaṃ nāma pañca kāmaguṇā, taṇhā vanatho \ldots Vanaṃ nāma nimittaggāho `itthī\,''ti vā `puriso\,''ti vā. Vanatho nāma tesaṃ tesaṃ aṅgapaccaṅgānaṃ anubyañjanaggāho \ldots Vanaṃ nāma cha ajjhattikabāhirāni āyatanāni apariññātāni. Yaṃ tadubhayaṃ paṭicca uppajjati saṃyojanaṃ, ayaṃ vanatho \ldots Vanaṃ nāma anusayo. Vanatho nāma pariyuṭṭhānaṃ.''} Four levels of meaning for ``forest'' (\textit{vana}) and ``underbrush'' (\textit{vanatha}): (1) sense-pleasure / craving; (2) general sign-grasping / specific feature-grasping; (3) uncomprehended sense-bases / fetters arising from them; (4) underlying tendency / active obsession. The Integration Mode holds all four readings simultaneously. This is the reference to ``cut down both the forest and the underbrush'' (Dhp 283).}

Four readings. All true. All pointing to the same instruction: cut the cause, and the effect falls with it.

\section*{Integration by Synonym}

The second type of integration brings together different formulations that point to the same thing. Liberation of the heart through the fading of lust and liberation through wisdom by the fading of ignorance --- these are synonyms for the trainee's fruit and the non-trainee's fruit, respectively. Or: the first is the fruit of the non-returner; the second is the highest fruit, arahantship. Or: the first is transcending the sense-desire realm; the second is transcending all three realms.\footnote{Pali: \textit{``Rāgavirāgā cetovimutti sekkhaphalaṃ; avijjāvirāgā paññāvimutti asekkhaphalaṃ. Idaṃ vevacanaṃ. Rāgavirāgā cetovimutti anāgāmiphalaṃ; avijjāvirāgā paññāvimutti aggaphalaṃ arahattaṃ. Idaṃ vevacanaṃ. Rāgavirāgā cetovimutti kāmadhātusamatikkamanaṃ; avijjāvirāgā paññāvimutti tedhātusamatikkamanaṃ. Idaṃ vevacanaṃ.''} Three synonym-identifications: (1) liberation of heart = trainee's fruit / liberation through wisdom = non-trainee's fruit; (2) liberation of heart = non-returner fruit / liberation through wisdom = arahantship; (3) liberation of heart = transcending the sense-desire realm / liberation through wisdom = transcending all three realms. Each pair is a different way of saying the same thing.}

And then the Netti brings in the long list of wisdom-synonyms we have already met: the wisdom-faculty, the wisdom-power, the higher wisdom training, the wisdom-aggregate, the investigation awakening-factor, equanimity, knowledge, right view, discernment, scrutiny, conscience, insight, knowledge of the teaching --- all synonyms, all pointing to the same capacity.\footnote{Pali: \textit{``Paññindriyaṃ, paññābalaṃ, adhipaññāsikkhā, paññākkhandho, dhammavicayasambojjhaṅgo, upekkhāsambojjhaṅgo, ñāṇaṃ, sammādiṭṭhi, tīraṇā, santīraṇā, hirī, vipassanā, dhamme ñāṇaṃ, sabbaṃ, idaṃ vevacanaṃ. Ayaṃ vevacanena samāropanā.''} Integration by synonym: a list of fourteen terms, all synonyms for wisdom (\textit{paññā}). The Integration Mode weaves them together so that wherever you encounter any one of these terms in a sutta, you recognize it as the same underlying capacity.}

This is the Synonym Mode's work folded into the Integration Mode: recognizing that different words in different suttas are pointing to the same reality, and reading accordingly.

\section*{Integration by Development}

The third type of integration shows how the various sets of practice-qualities develop together. The Buddha said: ``Dwell contemplating the body in the body, ardent, clearly comprehending, mindful, having removed covetousness and grief concerning the world.'' The Integration Mode maps each word to a faculty: ``ardent'' is the energy-faculty. ``Clearly comprehending'' is the wisdom-faculty. ``Mindful'' is the mindfulness-faculty. ``Having removed covetousness and grief'' is the concentration-faculty.\footnote{Pali: \textit{``Yathāha bhagavā `tasmātiha tvaṃ bhikkhu kāye kāyānupassī viharāhi, ātāpī sampajāno satimā vineyya loke abhijjhādomanassaṃ\,'. Ātāpīti vīriyindriyaṃ. Sampajānoti paññindriyaṃ. Satimāti satindriyaṃ. Vineyya loke abhijjhādomanassanti samādhindriyaṃ.''} The four qualities mentioned in the standard mindfulness formula, mapped to four faculties: ardent (\textit{ātāpī}) = energy-faculty; clearly comprehending (\textit{sampajāna}) = wisdom-faculty; mindful (\textit{satimā}) = mindfulness-faculty; having removed covetousness-grief (\textit{vineyya abhijjhādomanassa}) = concentration-faculty.}

And then: when you develop body-contemplation with these four faculties, the four foundations of mindfulness reach fulfillment. When the four foundations of mindfulness reach fulfillment, the four right strivings reach fulfillment. When the four right strivings reach fulfillment, the four bases of power reach fulfillment. When those reach fulfillment, the five faculties reach fulfillment. And so on through the entire set of awakening-qualities.\footnote{Pali: \textit{``Catūsu satipaṭṭhānesu bhāviyamānesu cattāro sammappadhānā bhāvanāpāripūriṃ gacchanti. Catūsu sammappadhānesu bhāviyamānesu cattāro iddhipādā bhāvanāpāripūriṃ gacchanti. Catūsu iddhipādesu bhāviyamānesu pañcindriyāni bhāvanāpāripūriṃ gacchanti. Evaṃ sabbe.''} The cascading development: four foundations of mindfulness → four right strivings → four bases of power → five faculties → and so on through all the qualities conducive to awakening. Each set reaching fulfillment causes the next to reach fulfillment.}

Why? Because all these awakening-qualities share a single characteristic: the characteristic of leading out. They all point in the same direction --- toward liberation. And because they share one characteristic, developing one develops them all.\footnote{Pali: \textit{``Kena kāraṇena? Sabbe hi bodhaṅgamā dhammā bodhipakkhiyā niyyānikalakkhaṇena ekalakkhaṇā, te ekalakkhaṇattā bhāvanāpāripūriṃ gacchanti. Ayaṃ bhāvanāya samāropanā.''} The key insight of integration by development: all awakening-related qualities (\textit{bodhipakkhiyā dhammā}) share the single characteristic of leading out (\textit{niyyānikalakkhaṇa}). Because they are ``of one characteristic'' (\textit{ekalakkhaṇa}), developing any one of them brings all the others to fulfillment. This is why the path is not a checklist of separate tasks but a single, integrated movement.}

This is why the path is not a checklist. It is a single movement that, when set in motion, carries everything forward.

\section*{Integration by Abandonment}

The fourth and final type of integration is the most detailed. It maps each of the four foundations of mindfulness to a specific network of abandonments:

Contemplating the body: you abandon the distortion that sees beauty in the unattractive. You comprehend physical food. You become free from clinging to sense-pleasure, free from the sense-pleasure yoke, detached from the body-tie of covetousness, free from the sense-pleasure taint. You cross the flood of sense-pleasure, remove the dart of lust, comprehend the form-based consciousness-station, abandon lust in the form-realm, and stop going by the wrong course of desire.\footnote{Pali: \textit{``Kāye kāyānupassī viharanto `asubhe subha\,''ti vipallāsaṃ pajahati, kabaḷīkāro cassa āhāro pariññaṃ gacchati, kāmupādānena ca anupādāno bhavati, kāmayogena ca visaṃyutto bhavati, abhijjhākāyaganthena ca vippayujjati, kāmāsavena ca anāsavo bhavati, kāmoghañca uttiṇṇo bhavati, rāgasallena ca visallo bhavati, rūpūpikā cassa viññāṇaṭṭhiti pariññaṃ gacchati, rūpadhātuyaṃ cassa rāgo pahīno bhavati, na ca chandāgatiṃ gacchati.''} Eleven abandonments from body-contemplation alone: (1) beauty-in-the-unattractive distortion, (2) physical food fully comprehended, (3) sense-pleasure clinging, (4) sense-pleasure yoke, (5) covetousness body-tie, (6) sense-pleasure taint, (7) sense-pleasure flood, (8) lust-dart, (9) form-based consciousness-station, (10) lust in the form-realm, (11) the wrong course of desire (\textit{chandāgati}).}

Contemplating feelings: you abandon the distortion that sees pleasure in suffering. You comprehend the nutriment of contact. You become free from clinging to existence, from the existence-yoke, from the body-tie of ill-will, from the existence-taint. You cross the flood of existence, remove the dart of hatred, comprehend the feeling-based consciousness-station, abandon lust in the feeling-realm, and stop going by the wrong course of hatred.\footnote{Pali: \textit{``Vedanāsu vedanānupassī viharanto `dukkhe sukha\,''ti vipallāsaṃ pajahati, phasso cassa āhāro pariññaṃ gacchati, bhavūpādānena ca anupādāno bhavati \ldots na ca dosāgatiṃ gacchati.''} Eleven parallel abandonments from feeling-contemplation, mapping onto the existence-related defilements rather than the sense-pleasure-related ones.}

Contemplating the mind: you abandon the distortion that sees permanence in the impermanent. You comprehend the nutriment of consciousness. You become free from clinging to views, from the view-yoke, from the body-tie of attachment to rules and rituals, from the view-taint. You cross the flood of views, remove the dart of conceit, comprehend the perception-based consciousness-station, abandon lust in the perception-realm, and stop going by the wrong course of fear.\footnote{Pali: \textit{``Citte cittānupassī viharanto `anicce nicca\,''ti vipallāsaṃ pajahati, viññāṇaṃ cassa āhāro pariññaṃ gacchati, diṭṭhupādānena ca anupādāno bhavati \ldots na ca bhayāgatiṃ gacchati.''} Eleven abandonments from mind-contemplation, now mapping onto the view-related defilements. Note the precise correspondence: each foundation of mindfulness maps onto one of the four distortions, one of the four nutrients, one of the four clingings, one of the four yokes, one of the four body-ties, one of the four taints, one of the four floods, one of the four darts, one of the four consciousness-stations, one of the four realm-based lusts, and one of the four wrong courses.}

Contemplating mental qualities: you abandon the distortion that sees self in what is not-self. You comprehend the nutriment of mental intention. You become free from clinging to self-doctrine, from the ignorance-yoke, from the body-tie of ``this alone is true,'' from the ignorance-taint. You cross the flood of ignorance, remove the dart of delusion, comprehend the formation-based consciousness-station, abandon lust in the formation-realm, and stop going by the wrong course of delusion.\footnote{Pali: \textit{``Dhammesu dhammānupassī viharanto `anattani attā\,''ti vipallāsaṃ pajahati, manosañcetanā cassa āhāro pariññaṃ gacchati, attavādupādānena ca anupādāno bhavati \ldots na ca mohāgatiṃ gacchati. Ayaṃ pahānena samāropanā.''} Eleven abandonments from dhamma-contemplation, mapping onto the delusion-related defilements. The four foundations of mindfulness, integrated by abandonment, yield forty-four specific abandonments in total (11 $\times$ 4). Each foundation clears a different sector of the mind's defilements. Together, they clear everything.}

\ornament

Four foundations. Four distortions. Four nutrients. Four clingings. Four yokes. Four body-ties. Four taints. Four floods. Four darts. Four consciousness-stations. Four realm-based lusts. Four wrong courses. Everything mapped to everything, everything integrated into one practice.

That is the Integration Mode. It does not add new content. It takes everything the other fifteen modes have revealed and shows that it all fits together --- that the Buddha's teaching is not a collection of separate insights but a single, unified system, visible from any angle, accessible through any doorway, and pointing always to the same freedom.\footnote{Pali: \textit{``Tenāha āyasmā mahākaccāyano --- `ye dhammā yaṃ mūlā, ye cekatthā pakāsitā muninā; te samāropayitabbā, esa samāropano hāro\,''ti.''} ``Those teachings that share a root, and those proclaimed by the Sage to have a single purpose --- they should be integrated. This is the Integration Mode.'' With this, the detailed analysis of the sixteen modes (\textit{hāravibhaṅga}) is complete. \textit{``Niṭṭhito ca hāravibhaṅgo.''}}
